\documentclass[UTF8,aspectratio=169]{ctexbeamer}
\usetheme{Madrid}
\usecolortheme{default}
\usepackage{graphicx}
\usepackage{booktabs}
\usepackage{amsmath}
\usepackage{hyperref}
\graphicspath{{figures/}{figures/beamer_pairs/}}

\setbeamertemplate{navigation symbols}{}
\setbeamerfont{frametitle}{size=\large}
\setbeamerfont{footnote}{size=\tiny}
\setbeamerfont{normal text}{size=\small}

\title{3D Gaussian Splatting 室内场景重建}
\author[Guo, Han, Wang]{Jingyi Guo \quad PB22011958 \\ Siqi Han \quad PB22010386 \\ Yiru Wang \quad PB22081510}
\date{2026 年 6 月}

\newcommand{\refline}[1]{\vfill{\tiny #1}}

\begin{document}

\begin{frame}
\titlepage
\end{frame}

\begin{frame}{任务要求}
\small
\begin{itemize}
\item 在共享室内测试场景上完成 3DGS 重建。
\item 输出可渲染的新视角图像，并计算 PSNR 和 SSIM。
\item 完成尺度标定和几何精度评估，给出多组测量结果。
\item 分析失败区域，给出有针对性的改进。
\item 提交技术文档、答辩材料、脚本、结果图、指标和模型文件。
\end{itemize}
\end{frame}

\begin{frame}{完整流程}
\small
有了共享图像数据，首先需要恢复相机从哪里拍摄。获得相机位姿后，三维高斯才能从正确视角渲染并与真实图像比较。
\begin{figure}
\centering
\includegraphics[width=0.94\linewidth]{pipeline.png}
\end{figure}
\end{frame}

\begin{frame}{COLMAP 和 SfM}
\small
有了多视角图像，需要估计相机位姿和稀疏三维点。COLMAP 使用 SfM 完成这一部分。
\[
\tilde{\mathbf{x}}_{ij} \sim K_i [R_i \mid \mathbf{t}_i]\tilde{\mathbf{X}}_j
\]
\begin{itemize}
\item SIFT 提取稳定特征点。
\item 特征匹配寻找不同图像中的同一物理点。
\item RANSAC 去除错误匹配。
\item 三角化恢复稀疏三维点。
\item Bundle Adjustment 同时优化相机和三维点。
\end{itemize}
\refline{Schönberger and Frahm 2016. Lowe 2004. Fischler and Bolles 1981. Triggs et al. 2000.}
\end{frame}

\begin{frame}{3DGS 表示}
\small
有了相机位姿和稀疏点云，下一步把稀疏点扩展为可训练的三维高斯。
\[
G_i(\mathbf{x}) =
\exp\left(-\frac{1}{2}
(\mathbf{x}-\boldsymbol{\mu}_i)^T
\Sigma_i^{-1}
(\mathbf{x}-\boldsymbol{\mu}_i)\right)
\]
\[
\Sigma_i = R_i S_i S_i^T R_i^T
\]
\begin{itemize}
\item 每个高斯包含位置、尺度、旋转、颜色和不透明度。
\item 高斯投影到图像平面后形成二维椭圆 footprint。
\item 训练会调整高斯参数，使渲染图接近真实图。
\end{itemize}
\refline{Kerbl et al. 2023.}
\end{frame}

\begin{frame}{可微渲染}
\small
有了三维高斯，模型需要从指定相机视角合成二维图像。多个高斯对同一像素贡献颜色时使用 alpha 合成。
\[
\mathbf{C} =
\sum_{i=1}^{N} T_i \alpha_i \mathbf{c}_i,
\quad
T_i = \prod_{k<i}(1-\alpha_k)
\]
\[
\mathcal{L} =
(1-\lambda)\|\hat{\mathbf{I}}-\mathbf{I}\|_1
+ \lambda(1-\mathrm{SSIM})
\]
\begin{itemize}
\item 可微渲染把图像误差反向传回三维高斯参数。
\item 本实验加入低学习率 finetune、有限增密和边缘一致性约束。
\end{itemize}
\refline{Kerbl et al. 2023.}
\end{frame}

\begin{frame}{Mip Splatting：为什么需要它}
\small
普通 3DGS 在投影后使用固定的屏幕空间 dilation：
\[
G^{2D}(\mathbf{x}) =
\exp\left(
-\frac{1}{2}
(\mathbf{x}-\boldsymbol{\mu})^T
(\Sigma^{2D}+s\,\mathbf{I})^{-1}
(\mathbf{x}-\boldsymbol{\mu})
\right)
\]
\begin{itemize}
\item 固定 \(s\mathbf{I}\) 只在训练采样率附近近似合理。
\item 视角拉远或焦距变小，dilation 会把能量铺到过大的像素区域，产生变亮和膨胀。
\item 视角拉近或焦距变大，训练得到的过小三维高斯会暴露为高频伪影和边缘侵蚀。
\item Mip Splatting 的核心是先限制三维表示频率，再用二维 Mip filter 模拟像素面积积分。
\end{itemize}
\refline{Yu et al. 2024.}
\end{frame}

\begin{frame}{Mip Splatting：三维平滑滤波}
\scriptsize
对第 \(k\) 个高斯，先用训练相机估计它在三维空间的最高可恢复采样率。若该高斯到第 \(n\) 个可见相机的深度为 \(d_{n,k}\)，焦距为 \(f_n\)，则
\[
\hat{T}_{n,k}=\frac{d_{n,k}}{f_n},
\quad
\hat{\nu}_k =
\max_{n}
\left(
\mathbf{1}_{n}(\boldsymbol{\mu}_k)
\frac{f_n}{d_{n,k}}
\right)
\]
\[
G^{reg}_k(\mathbf{x})
=
(G_k \otimes G_{low})(\mathbf{x}),
\quad
\Sigma^{reg}_k =
\Sigma_k+\frac{s_3}{\hat{\nu}_k^2}\mathbf{I}
\]
\[
\alpha^{reg}_k =
\alpha_k
\sqrt{
\frac{|\Sigma_k|}
{|\Sigma_k+\frac{s_3}{\hat{\nu}_k^2}\mathbf{I}|}
}
\]
\begin{itemize}
\item \(\hat{\nu}_k\) 取所有可见训练视角中的最大采样率，等价于使用最细的可靠观察尺度。
\item 与高斯低通卷积只需把协方差相加，代价低，可在训练中周期性更新。
\item 行列式系数补偿体积变化，避免平滑后 opacity 人为变亮。
\end{itemize}
\refline{Yu et al. 2024.}
\end{frame}

\begin{frame}{Mip Splatting：二维 Mip 滤波}
\scriptsize
三维滤波解决重建本身的过高频率；渲染到低采样率视角时，还需要在图像平面模拟一个像素面积内的积分。
\[
\Sigma^{2D}_k =
J_k R \Sigma^{reg}_k R^T J_k^T
\]
\[
G^{2D,mip}_k(\mathbf{x}) =
\sqrt{
\frac{|\Sigma^{2D}_k|}
{|\Sigma^{2D}_k+s_2\mathbf{I}|}
}
\exp\left(
-\frac{1}{2}
(\mathbf{x}-\boldsymbol{\mu}^{2D}_k)^T
(\Sigma^{2D}_k+s_2\mathbf{I})^{-1}
(\mathbf{x}-\boldsymbol{\mu}^{2D}_k)
\right)
\]
\begin{itemize}
\item 普通 dilation 只放大 footprint，不降低峰值；Mip filter 同时放大协方差并衰减峰值。
\item \(s_2\mathbf{I}\) 近似一个像素的 box filter，使单像素内能量守恒，更符合成像过程。
\item 本实验在 Sequence 04 上训练 Mip Splatting 12000 次，达到 PSNR 27.77、SSIM 0.811，接近最终 3DGS finetune。
\end{itemize}
\refline{Yu et al. 2024.}
\end{frame}

\begin{frame}{PSNR 和 SSIM}
\small
在汇报结果之前，先定义两个图像质量指标。
\[
\mathrm{MSE}=
\frac{1}{HW}\sum_{u,v}
\|\hat{\mathbf{I}}(u,v)-\mathbf{I}(u,v)\|_2^2
\]
\[
\mathrm{PSNR}=10\log_{10}\frac{1}{\mathrm{MSE}}
\]
\[
\mathrm{SSIM}(x,y)=
\frac{(2\mu_x\mu_y+C_1)(2\sigma_{xy}+C_2)}
{(\mu_x^2+\mu_y^2+C_1)(\sigma_x^2+\sigma_y^2+C_2)}
\]
\begin{itemize}
\item PSNR 越高，像素误差越小。
\item SSIM 越接近 1，结构相似度越高。
\end{itemize}
\refline{Wang et al. 2004.}
\end{frame}

\begin{frame}{COLMAP 结果}
\small
\begin{tabular}{lrrrr}
\toprule
序列 & 输入 & 注册 & 注册率 & 稀疏点 \\
\midrule
Sequence 01 & 187 & 187 & 1.000 & 41327 \\
Sequence 02 & 80 & 49 & 0.613 & 13667 \\
Sequence 03 & 120 & 120 & 1.000 & 54290 \\
Sequence 04 & 80 & 80 & 1.000 & 40064 \\
Sequence 05 & 422 & 422 & 1.000 & 135204 \\
\bottomrule
\end{tabular}
\vspace{0.4cm}

除 Sequence 02 外，其余序列均完整注册。Sequence 02 注册率较低，说明该片段存在弱纹理或匹配不足。
\end{frame}

\begin{frame}{调参策略}
\small
\begin{itemize}
\item 数据侧采用清晰度筛选和局部分区，减少模糊帧和弱覆盖区域影响。
\item COLMAP 侧检查注册率和稀疏点云覆盖情况。
\item 3DGS 侧采用较长训练、低学习率 finetune、有限增密和边缘一致性项。
\item 进阶方法侧实跑 Mip Splatting、2DGS 和 Depth RegularizedGS。
\end{itemize}
\begin{tabular}{lrrr}
\toprule
序列 & 基线 PSNR & 最终 PSNR & 增益 \\
\midrule
Sequence 03 & 23.42 & 26.33 & 2.91 \\
Sequence 04 & 25.30 & 28.36 & 3.05 \\
Sequence 05 & 25.17 & 26.39 & 1.21 \\
\bottomrule
\end{tabular}
\end{frame}

\begin{frame}{最终渲染指标}
\small
\begin{tabular}{lrrrr}
\toprule
序列 & 图像数 & 高斯数 & PSNR & SSIM \\
\midrule
Sequence 01 & 187 & 69507 & 26.36 & 0.836 \\
Sequence 02 & 49 & 20718 & 30.15 & 0.884 \\
Sequence 03 & 120 & 64675 & 26.33 & 0.813 \\
Sequence 04 & 80 & 50663 & 28.36 & 0.811 \\
Sequence 05 & 422 & 145975 & 26.39 & 0.812 \\
\bottomrule
\end{tabular}
\begin{figure}
\centering
\includegraphics[width=0.58\linewidth]{final_render_quality.png}
\end{figure}
\end{frame}

\begin{frame}{Sequence 01 结果}
\begin{figure}
\centering
\includegraphics[height=0.68\textheight,keepaspectratio]{sequence_01_pair_1.jpg}
\end{figure}
\end{frame}

\begin{frame}{Sequence 02 结果}
\begin{figure}
\centering
\includegraphics[height=0.68\textheight,keepaspectratio]{sequence_02_pair_1.jpg}
\end{figure}
\end{frame}

\begin{frame}{Sequence 03 结果}
\begin{figure}
\centering
\includegraphics[height=0.68\textheight,keepaspectratio]{sequence_03_pair_1.jpg}
\end{figure}
\end{frame}

\begin{frame}{Sequence 04 结果}
\begin{figure}
\centering
\includegraphics[height=0.68\textheight,keepaspectratio]{sequence_04_pair_1.jpg}
\end{figure}
\end{frame}

\begin{frame}{Sequence 05 结果}
\begin{figure}
\centering
\includegraphics[height=0.68\textheight,keepaspectratio]{sequence_05_pair_1.jpg}
\end{figure}
\end{frame}

\begin{frame}{进阶方法对比}
\scriptsize
\begin{tabular}{lrr}
\toprule
方法 & PSNR & SSIM \\
\midrule
受控 3DGS finetune & 28.36 & 0.811 \\
Mip Splatting & 27.77 & 0.811 \\
2DGS & 27.18 & 0.796 \\
Depth RegularizedGS & 15.92 & 0.662 \\
\bottomrule
\end{tabular}
\vspace{0.05cm}
\begin{center}
\includegraphics[height=0.30\textheight,keepaspectratio]{advanced_methods_comparison.png}
\end{center}
\refline{Yu et al. 2024. Huang et al. 2024. Chung et al. 2024.}
\end{frame}

\begin{frame}{Mip Splatting 结果}
\begin{figure}
\centering
\includegraphics[height=0.66\textheight,keepaspectratio]{sequence04_mip_pair.jpg}
\end{figure}
\refline{Yu et al. 2024.}
\end{frame}

\begin{frame}{2DGS 结果}
\begin{figure}
\centering
\includegraphics[height=0.66\textheight,keepaspectratio]{sequence04_twodgs_pair.jpg}
\end{figure}
\refline{Huang et al. 2024.}
\end{frame}

\begin{frame}{Depth RegularizedGS 结果}
\begin{figure}
\centering
\includegraphics[height=0.66\textheight,keepaspectratio]{sequence04_depth_regularized_pair.jpg}
\end{figure}
\refline{Chung et al. 2024.}
\end{frame}

\begin{frame}{几何精度}
\small
\[
\mathbf{q}_i \approx sR\mathbf{p}_i+\mathbf{t}
\]
\begin{tabular}{lrrrr}
\toprule
序列 & 平均误差 m & 中位误差 m & 20 cm 内 & 10 cm 内 \\
\midrule
Sequence 01 & 0.277 & 0.195 & 0.510 & 0.254 \\
Sequence 02 & 0.740 & 0.256 & 0.433 & 0.235 \\
Sequence 03 & 0.280 & 0.204 & 0.494 & 0.267 \\
Sequence 04 & 0.246 & 0.155 & 0.602 & 0.327 \\
Sequence 05 & 0.251 & 0.188 & 0.523 & 0.286 \\
\bottomrule
\end{tabular}
\end{frame}

\begin{frame}{失败区域}
\small
\begin{itemize}
\item 白色弱纹理柜面和墙面缺少稳定梯度，COLMAP 特征和图像监督都较弱。
\item 反光区域存在视角相关高光，静态颜色高斯难以同时解释多个视角。
\item 黑色线条、门框、线缆、文字和地毯色块恢复较稳定，因为它们提供高对比特征。
\item 后续可加入反光 mask、可靠深度正则、线段监督、平面约束、Mip Splatting 和分区训练。
\end{itemize}
\end{frame}

\begin{frame}{结论}
\small
\begin{itemize}
\item 五个序列的平均 PSNR 均超过 25 dB。
\item Sequence 04 达到 28.36 dB，满足优秀档位。
\item Sequence 04 和 Sequence 05 的全局中位几何误差低于 20 cm，局部区域达到 10 cm 附近。
\item Mip Splatting 和 2DGS 均完成实跑，指标接近最终 3DGS finetune。
\item Depth RegularizedGS 完成配置、训练和渲染，但在当前稀疏深度设置下效果较弱。
\end{itemize}
\end{frame}

\begin{frame}{交付内容}
\small
\begin{itemize}
\item 完整实验报告 PDF。
\item 答辩 beamer PDF。
\item 五个序列的渲染图、真实图对比和指标表。
\item 几何评估脚本和精度数据。
\item 可被 SuperSplat 或 3DGS viewer 打开的 PLY 模型。
\item Mip Splatting、2DGS、Depth RegularizedGS 的训练输出和对比结果。
\item 可复现实验的全部脚本和成员信息。
\end{itemize}
\end{frame}

\end{document}
